\chapter{Square-Root Sampling Branch}\label{cap:bagoftricks}

In this Chapter, we propose a simple and effective framework called
Square-Root Sampling Branch (SSB) to improve long-tail visual recognition of
animal species in camera-trap images. We systematically evaluate it against
state-of-the-art methods for handling this task: square-root sampling,
class-balanced focal loss, and balanced group softmax.

\section{Introduction}

One major problem with camera-trap data is the extreme imbalance among
categories as shown in \autoref{fig:train_dist}, in which some species have many
images while others have just a few, following a long-tail distribution. As
reinforced by \citet{schneider2020three}, models generally achieve poor recall
scores for species with fewer training images, although these rare species are
often the ones of greatest interest to ecologists.


Many works in the literature on machine learning have proposed techniques for
addressing the long-tailed issue, such as re-sampling the data
\citep{Kang2020Decoupling}, re-weighting the loss \citep{cui2019class}, or
creating multiple branches in the networks \citep{li2020overcoming,
zhou2020bbn, cai2021ace}. Despite this, most of the works for classifying animal
species in camera-trap images do not apply any method for handling long-tail
classification. Some works remove the tail classes \citep{villa2017towards,
willi2019identifying} or group them into generic
categories~\citep{zualkernan2020towards}.

Based on the empirical observation that square-root
sampling~\citep{Kang2020Decoupling} was the method that most improved the
performance of the tail classes in our preliminary experiments, though at the cost of strongly degrading the
head classes, in this chapter a framework called Square-root Sampling Branch
(SSB) is proposed, which combines the predictions of two classification
branches: one trained using square-root sampling and the other one trained
without any additional method for long-tail classification.

As the contributions, we can highlight:
\begin{itemize}
\item The proposal of a simple and effective framework called Square-root
Sampling Branch (SSB) for long-tail visual recognition, which achieves a better
trade-off as it improves performance in tail classes at a low cost of head
classes’ accuracy when compared to the state-of-the-art methods.

\item New baselines in camera-trap datasets are established using training/test
partitioning that is based on locations by combining the most recent training
tricks with both traditional architectures, such as ResNet-50, and newer ones
such as Swin Transformer.

\item Our experiments show that Swin Transformer can achieve excellent results
in camera-trap datasets even without using any additional method for handling
the long-tail issue.
\end{itemize}

\section{Materials and Methods}

Our methodology was built to assess the effectiveness of our proposed SSB
framework in improving the recognition performance in the tail classes and
comparing it with methods available in the literature. First, a strong baseline
was trained with several recent training tricks available in the literature.
Then, each long-tail method was applied and evaluated using appropriate metrics
to assess long-tail performance. To reach a more general conclusion, the
performance of the methods on four long-tail camera-trap datasets with different
characteristics and four families of computer vision models were evaluated.

In this section, the datasets, the chosen computer vision architectures,
and the baseline training procedure are described. Then, the families of methods
for long-tail recognition and the selected methods used in this work are
presented. Finally, a description of our proposed SSB framework is given.

\subsection{Datasets}

This subsection describes the four camera-trap datasets used in this study:
Snapshot Serengeti, Caltech Camera Traps, WCS Camera Traps, and Wellington
Camera Traps. These datasets were obtained from different regions around the
globe and have different characteristics, such as the number of training
instances, imbalance factor, and available species. While annotations for these
datasets are provided for bursts (i.e., a sequence of images), we assign
each image the label inherited from its burst annotations during the training
process. However, this approach may result in mislabeling, especially when the
animal appears in only a few images within a burst. Since manual verification
was not feasible due to the large amount of images in the datasets, the data
was used as provided, as has been done in prior studies \citep{villa2017towards,
norouzzadeh2018automatically}. 

Nonetheless, it is important to note that the
recommended approach is to have annotations at an image-level, as this provides
more accurate training data. Following the recent recommendations, the datasets
were split into training/validation/test partitions based on camera locations to
better assess the models’ generalization. Additionally, images were removed from
all classes which, due to the location-based split, do not have instances in the
training or test partitions. Images containing more than one species as well as
corrupted pictures were also removed. \autoref{table:datasets_instances}
summarizes the number of images used from each dataset in this work. All of
these datasets follow a long-tail distribution, as shown in
\autoref{fig:train_dist}, in which a few classes (head) concentrate most of the
instances and the remaining classes have just a few samples in comparison.

\begin{table}[htb]
\resizebox{\textwidth}{!}{
\begin{threeparttable}
\begin{center}
\caption{Number of images and locations per partition used from each dataset.}
\label{table:datasets_instances}
\begin{tabular}{llllllll}
\hline
 &  & \multicolumn{2}{c}{Training} &
\multicolumn{2}{c}{Validation} & \multicolumn{2}{c}{Test} \\ \cline{3-8}
Dataset & Classes & \# loc  & \# instances & \# loc & \# instances

  & \# loc   & \# instances  \\ \hline
Snapshot Serengeti & 47 & 156 & 1,059,793  & 23\textsuperscript{*}


  & 360,317   & 46\textsuperscript{*}  & 738,404
 \\
Caltech CT & 19 & 80  & 122,693   & 20\textsuperscript{*}


  & 57,162
  & 40\textsuperscript{*}  & 61,665         \\
WCS & 247 & 2,506  & 599,040   & 313  & 79,906


  & 314 & 62,734         \\
Wellington CT & 14 & 111 & 160,937   & 36   & 49,127


  & 35  & 55,424         \\ \hline
\end{tabular}
\begin{tablenotes}
\footnotesize
\item [*] As long as there are no official test sets for Snapshot Serengeti and
Caltech CT datasets, their recommended validation sets were used as the test
sets in this work. For these datasets, validation sets (which were called
minival) were held out from their respective training sets following the
locations selected in Chapter~\ref{cap:filtering}.
\end{tablenotes}
\end{center}
\end{threeparttable}}
\end{table}

\paragraph{\textbf{Snapshot Serengeti}~\citep{dryad_5pt92}} This dataset is
composed of more than 7 million images obtained from 225 locations spread
across the Serengeti National Park in Tanzania, and is divided into 11 seasons.
To maintain consistency with the literature, only the first six seasons are used
in this work. The recommended training/validation partitions based on locations
were used \citep{ss_lila}, with the validation set being used to evaluate the
models. To adjust the hyperparameters, a mini-validation (minival) partition
was held out from the training set using the locations selected in
Chapter~\ref{cap:filtering}. Specifically for this dataset, given that images
from empty bursts (background or blank) are near identical, only one picture was
picked at random from each empty burst.  This also helps to keep a tractable
number of images to train the models. In total, approximately 2 million pictures
distributed in 47 categories were used.

\paragraph{\textbf{Caltech Camera Traps}~\citep{beery2018recognition}} This
dataset is composed of approximately of 240,000 images from 140 locations in
the southwestern United States. The recommended location-based
training/validation partitions were used \citep{caltech_lila}, and, as for the
Snapshot Serengeti, the validation set was used for testing, and a
mini-validation partition was created from the training set for hyperparameter
adjustment based on locations selected in
Chapter~\ref{cap:filtering}.

\paragraph{\textbf{WCS Camera Traps}~\citep{wcs_lila}} This dataset is provided
by the Wildlife Conservation Society, and it is composed of 1.4 million images
of around 675 species from 12 countries, thus making it one of the most diverse
datasets publicly available. In our experiments, the recommended
training/validation/test partitions based on locations were used. However,
since some locations were not allocated to any partition, images from them were
not used. The generic unknown categories and those not related to animals, such
as misfire, start, and end were also manually removed. As images from classes
without instances in training or testing sets were also removed, the resulting
selected subset is composed of images from 247 categories.

\paragraph{\textbf{Wellington Camera Traps}~\citep{anton2018}} This dataset is
originally composed of around 270,000 images collected from 187 camera locations
in the Wellington region of New Zealand. The training/validation/test proposed
by \citet{shashidhara2020sequence} was used. Images marked as unclassifiable
were removed, in addition to images from the categories not present in the
training or test partitions, thus resulting in a dataset of around 265,000
images from 14 categories. Unlike the other datasets, where the empty class is 
the majority category, in this dataset the empty class has the second-largest 
number of instances -- with the generic bird category appearing three times more 
frequently.

\subsection{Architectures}

Four popular computer vision architectures with different purposes and
characteristics were chosen to assess whether long-tail classification methods
work well on different backbones: ResNet-50 \citep{he2016deep},
MobileNetV3-Large \citep{howard2019searching},
EfficientNetV2-B2 \citep{tan2021efficientnetv2}, and Swin-S
\citep{liu2021swin}.
First, ResNet-50 was selected because it is the most common baseline for almost
any method in computer vision. As using computer vision models directly on edge
devices for processing camera-trap data has gained attention in the literature
\citep{elias2017s, zualkernan2020towards, Cunha2021, zualkernan2022iot},
MobileNetV3-Large was chosen as a benchmark for this situation. For the
state-of-the-art families of EfficientNet \citep{tan2019efficientnet} and Swin
Transformer \citep{liu2021swin} models, lightweight versions --
EfficientNetV2-B2 and Swin-S, respectively -- were chosen, since most of the
camera-trap practitioners usually do not have access to powerful computers to
run big models \citep{schneider2020three}.

\subsection{Baseline training settings}\label{sec:baseline_settins}

As shown in several works \citep{he2019bag, bello2021revisiting, wightman2021resnet},
improved training strategies are responsible for significant gains in the
accuracy of classical architectures such as ResNet-50. Following these ideas,
several tricks were included in our baseline training procedure, and these are
described below.

\paragraph{\textbf{Image preprocessing}} Initially, all images were resized so
that the largest dimension was at most 768 px, and the smallest one was at
least 450 px. During training, first, a rectangular region of aspect ratio
sampled in $[3/4, 4/3]$ and area in $[65\%, 100\%]$ was randomly cropped and
resized to the input size of the architecture. Next, the image was flipped
horizontally with a 50\% probability. Then,
RandAugment~\citep{cubuk2020randaugment} was used with parameters $N=2$ and
$M=9$. Finally, the image was rescaled to match the input scale of each
architecture, as shown in \autoref{table:arch_rescale}. During the evaluation,
the preprocessing consists only of resizing the image to the input architecture
size and rescaling.

\begin{table}[htb]
\caption{Input scaling procedure for each architecture.}
\label{table:arch_rescale}
\begin{center}
\begin{tabular}{lp{0.7\textwidth}}
\hline
Architecture & Scaling Procedure
                     \\ \hline
ResNet-50   & Convert from RGB to BGR, zero-center w.r.t. the ImageNet
dataset, without scaling \\
MobileNetV3-Large  & Pixels between 0 and 255 (model includes preprocessing
layer)                     \\
EfficientNetV2-B2 & Scale pixels between -1 and 1
                      \\
Swin-S  & Scale pixels between 0 and 1, normalize w.r.t. the ImageNet
dataset               \\ \hline
\end{tabular}
\end{center}
\end{table}

\paragraph{\textbf{Training procedure}} The models were initialized with
ImageNet pre-trained weights, replacing the classifier layer, and, then, were
trained with a batch size of 64 for 30 epochs when fine-tuning the whole model
or for 12 epochs when training only the classifier layer. The AdamW optimizer
\citep{loshchilov2018decoupled} was employed for training with an initial
learning rate of $10^{-5}$ and weight decay of $10^{-7}$. The learning rate was
linearly warmed up from $0$ to the initial value for two epochs when fine-tuning
-- one epoch when training only the classifier --, and then decayed to $0$
using the cosine schedule \citep{he2019bag}.

\subsection{Training approaches for long-tail visual recognition}

To achieve more general conclusions, we performed a systematic study using
state-of-the-art methods for handling the long-tail visual recognition problem.
As described in Section~\ref{sec:func_longtail}, we grouped them into four main
families according to their focus: re-sampling the data, re-weighting the loss,
using two stages to focus on representation learning, and creating multiple
branches in the network, selecting one method from each family to run our
experiments. \autoref{fig:train_diagram} illustrates a functional summary of
the model training process, highlighting the stages where each method studied in
this work operates.

For representing the re-weighting methods, we selected the square-root
resampling~\cite{Kang2020Decoupling, zhang2021bag} by setting $q=1/2$ in
Equation~\ref{eq:resampling}, as this approach returns a less imbalanced
dataset. For the re-weighting methods, we used the Class-Balanced Focal
(CB-Focal) loss, which is a combined version of Focal loss and Class-Balanced
loss. In our experiments, we set the hyperparameters $\beta$ to $0.9$ and
$\gamma$ to $2.0$. In our experiments, we also compare one-stage and two-stage
training strategies. In the one-stage approach, the balancing method is applied
end-to-end, while in the two-stage method~\cite{Kang2020Decoupling}, the model
is initially trained using the original unbalanced training distribution. Then,
only the classifier is trained using a balancing method, such as re-weighting or
re-sampling. Among the various multi-branch network approaches available in the
literature, we included BAGS~\cite{li2020overcoming} in our experiments.

The method of using an object detector to crop regions of interest from
camera-trap images before classifying them was also included in our experiments.
Despite not dealing directly with the imbalance problem, it has been shown in
the literature \citep{schneider2018deep, beery2019efficient} that this trick
improves the recognition score in camera-trap images. In our experiments, we
used MegaDetectorV4~\citep{beery2019efficient} to generate bounding box
predictions, considering only those with a confidence score higher than $0.6$;
otherwise, the entire image was used. In this operation, we applied a squared
crop around the bounding box with the highest confidence, executed before the
random crop operation in the preprocessing pipeline
(\autoref{sec:baseline_settins}).

\subsection{SSB: Square-root Sampling Branch}

Our proposed framework is a multi-branch network approach that was inspired by
an empirical observation in our preliminary experiments: the square-root
sampling strategy was the method that most improved the accuracy of the tail
classes over the baseline; however, at the cost of significantly decreasing the
performance of the head classes. The question to answer is: is it possible to
combine the \textbf{S}quare-root \textbf{S}ampling \textbf{B}ranch (SSB)
predictions with the baseline to improve the tail classes with minimal
performance degradation for the head classes?

As illustrated in \autoref{fig:ssb}, the proposed framework is trained in a
two-stage fashion. First, the classifier ($f_i$) and the backbone are trained
jointly using the original data distribution (instance-based sampling). For the
second stage, the network backbone is frozen and the classifier is replaced by
a new randomly initialized one. Then, this new classifier ($f_{sqrt}$) is
trained using data sampled by the square-root sampling strategy. Both
classifiers are connected on top of the backbone during the inference. Their
predictions are aggregated by masking the softmax outputs to keep only the head
classes probabilities from the $f_i$ head and the remaining classes
probabilities from the $f_{sqrt}$ head.

The $C$ classes were divided into four groups based on the number of training
instances, in which a class $j$ with $n_j$ instances belongs to group $G_k$ if
$s_k^l \leq n_j < s_k^h$. The limits of each group are set to $(0, 10)$, $(10,
10^2)$, $(10^2, 10^3)$, and $(10^3, +\infty)$. However, SSB does not have a
special group for the background class, with this class belonging to the group
$G_4$. Given a feature vector $h$ generated by the backbone and the predictions
$p_i = f_i(h)$ and $p_{sqrt} = f_{sqrt}(h)$ generated by each classification
head with softmax, the final remapped prediction vector $p_r$ is given by:
\begin{equation}\label{eq:ssb}
 p_r = f_r(p_i, p_{sqrt}) = Qp_i + (\bm{1} - Q)p_{sqrt},
\end{equation}
where $Q$ is a matrix $C\times C$ that masks the softmax outputs based on which
group a class belongs to:
\begin{equation}\label{eq:ssb_mask}
 Q = \left\{ q_{a,b} | q_{a, b} = \begin{cases}
      1 & \text{if } a=b \text{ and the class } a \in G_4 \text{,}\\
      0 & \text{otherwise.}
    \end{cases}\right\}.
\end{equation}
As with BAGS post-processing, the final prediction vector produced by SSB is
not
a real probability vector because its sum is not equal to 1.

\section{Experiments and Results}
Our experiments were broken down into three main series. In the first, our
proposed method SSB is compared to three two-stage methods for long-tail
problems, as well as the baseline. In the second series, the same methods are
compared, but using cropped images as inputs instead of the entire images.
Finally, in the third series, SSB is compared to single-stage training methods.
The results presented in this section were obtained using the test sets
specified for each dataset in \autoref{table:datasets_instances}.

\subsection{Comparison between SSB and two-stage methods for long-tail visual
recognition}

In the first experiment, the effectiveness of our proposed SSB method when
using the full image as input for the training and testing stages was evaluated.
As our proposed method is trained using a two-stage strategy, in this
experiment, only the results with two-stage training methods were compared.
Initially, each architecture was trained using our baseline training settings,
and then each long-tail method was applied to the classifier layer, though
keeping the backbone frozen during the second-stage training. This allows a fair
comparison between the methods in terms of the representation learned by the
backbone.

The accuracy results are shown in \autoref{table:fullimage_results}. The
proposed SSB consistently improved accuracy for $Bin_1$, $Bin_2$, and $Bin_3$
when compared to the baseline for almost all cases, with a minimal loss in the
head class accuracy ($Bin_4$), and achieved the best or the second-best
accuracy per bin for most of the cases. The best results for $Bin_1$, $Bin_2$,
and $Bin_3$ were achieved by the square-root sampling, with an improvement
ranging from 3\% to more than 15\% over the baseline. However, such an
improvement came at the cost of reducing the accuracy of the majority classes
($Bin_4$) by at least 3\%. This aspect may be observed in the overall accuracy,
since the square-root sampling achieved the worst results among all the methods.
On the other hand, BAGS achieved the best overall accuracy results for some
cases; while, in others, it even reduced the accuracy for the tail classes
compared to the baseline.  All methods, however, failed to improve the
performance on tail classes ($Bin_2$) for the smaller datasets (Caltech and
Wellington). In this case, accuracies were close to zero.

\begin{landscape}

\begin{table}[h]
\footnotesize
\caption[Comparison between SSB accuracy with different methods for long-tail
classification of different models and datasets when using the entire image as
the input]{Comparison between SSB accuracy with different methods for long-tail
classification of different models and datasets when using the entire image as
the input. Bold and underlined values, respectively, indicate the best and
second-best result for each model and dataset combination. $Acc_i$ is the top-1
accuracy for the $Bin_i$. As some category bins do not occur in some of the
datasets, in this case, the corresponding column was removed.}
\label{table:fullimage_results}
\resizebox{1.5\textwidth}{!}{
\begin{tabular}{ll|lllll|llll|lll|llll}
\hline
\multicolumn{1}{l}{}   & \multicolumn{1}{l|}{} &
\multicolumn{5}{c|}{\textbf{WCS}}




 & \multicolumn{4}{c|}{\textbf{Snapshot Serengeti}}


  &
\multicolumn{3}{c|}{\textbf{Caltech CT}}


  & \multicolumn{4}{c}{\textbf{Wellington CT}}


                                                               \\
\multicolumn{1}{l}{\textbf{Model}} & \multicolumn{1}{l|}{\textbf{Training}} &
\multicolumn{1}{c}{\textbf{$Acc_1$}} & \multicolumn{1}{c}{\textbf{$Acc_2$}} &
\multicolumn{1}{c}{\textbf{$Acc_3$}} & \multicolumn{1}{c}{\textbf{$Acc_4$}} &
\multicolumn{1}{c|}{\textbf{$Acc_{all}$}} &
\multicolumn{1}{c}{\textbf{$Acc_2$}} & \multicolumn{1}{c}{\textbf{$Acc_3$}} &
\multicolumn{1}{c}{\textbf{$Acc_4$}} &
\multicolumn{1}{c|}{\textbf{$Acc_{all}$}}&
\multicolumn{1}{c}{\textbf{$Acc_2$}} & \multicolumn{1}{c}{\textbf{$Acc_4$}} &
\multicolumn{1}{c|}{\textbf{$Acc_{all}$}} &
\multicolumn{1}{c}{\textbf{$Acc_2$}} & \multicolumn{1}{c}{\textbf{$Acc_3$}} &
\multicolumn{1}{c}{\textbf{$Acc_4$}} & \multicolumn{1}{c}{\textbf{$Acc_{all}$}}
\\ \hline
\multirow{5}{*}{\textbf{ResNet-50}} & \textbf{Baseline} & 0.00 & 13.05 & 32.06 &
\textbf{85.02} & {\ul 81.53} & 2.90 & 44.82 & {\ul 92.96} & {\ul 92.79} & 0.00 &
44.95 & 42.19 & 0.0 & 4.58 & \textbf{75.66} & 72.23 \\
 & \textbf{Sqrt-Samp} & \textbf{16.25} & \textbf{28.94} & \textbf{48.98} & 75.38
& 73.47 & {\ul 6.52} & \textbf{51.58} & 87.29 & 87.16 & 0.00 & 39.59 & 37.16 &
0.0 & {\ul 10.65} & 70.91 & 68.01 \\
 & \textbf{CBFocal} & 5.00 & 14.42 & 40.61 & 83.41 & 80.43 & 0.72 & 38.66 &
92.88 & 92.69 & 0.00 & 44.47 & 41.74 & 0.0 & 4.84 & 75.12 & 71.73 \\
 & \textbf{BAGS} & \textbf{16.25} & 14.19 & 34.84 & 84.43 & 81.14 &
\textbf{9.42} & 41.27 & \textbf{93.12} & \textbf{92.94} & \textbf{1.56} &
\textbf{48.04} & \textbf{45.19} & 0.0 & \textbf{25.92} & 75.00 & \textbf{72.64}
\\
 & \textbf{SSB (Ours)} & {\ul 15.00} & {\ul 21.79} & {\ul 40.97} & {\ul 84.85} &
\textbf{81.91} & 5.07 & {\ul 48.16} & 92.91 & 92.75 & {\ul 0.03} & {\ul 46.09} &
{\ul 43.26} & 0.0 & 5.51 & {\ul 75.62} & {\ul 72.24} \\ \hline
\multirow{5}{*}{\textbf{MBNetV3}} & \textbf{Baseline} & 0.00 & 11.92 & 36.52 &
\textbf{84.24} & {\ul 80.98} & 2.17 & 29.78 & \textbf{92.73} & {\ul 92.52} &
0.00 & \textbf{46.14} & \textbf{43.31} & 0.0 & 5.70 & \textbf{75.65} & {\ul
72.28} \\
 & \textbf{Sqrt-Samp} & {\ul 7.50} & \textbf{28.26} & \textbf{46.85} & 76.71 &
74.61 & {\ul 4.35} & \textbf{43.95} & 87.62 & 87.47 & \textbf{0.16} & 37.76 &
35.45 & 0.0 & {\ul 21.68} & 70.45 & 68.10 \\
 & \textbf{CBFocal} & 0.00 & 9.19 & 36.98 & 79.41 & 76.42 & 0.72 & 29.78 & 92.44
& 92.23 & 0.00 & 41.27 & 38.74 & 0.0 & 4.58 & 74.94 & 71.55 \\
 & \textbf{BAGS} & \textbf{10.00} & 10.44 & 28.68 & {\ul 84.09} & 80.48 &
\textbf{5.07} & 28.78 & 92.19 & 91.97 & 0.00 & 43.96 & 41.26 & 0.0 &
\textbf{23.41} & 74.02 & 71.58 \\
 & \textbf{SSB (Ours)} & 1.25 & {\ul 20.32} & {\ul 40.43} & 83.99 &
\textbf{81.04} &
3.62 & {\ul 36.97} & {\ul 92.72} & \textbf{92.53} & {\ul 0.11} & {\ul 45.93} &
{\ul 43.12} & 0.0 & 13.54 & {\ul 75.34} & \textbf{72.36} \\ \hline
\multirow{5}{*}{\textbf{EffV2-B2}} & \textbf{Baseline} & 0.00 & 10.22 & 32.31 &
86.35 & 82.75 & 2.90 & 53.71 & 93.82 & 93.68 & {\ul 0.08} & 56.36 & 52.91 & 0.0
& 5.70 & {\ul 77.11} & 73.67 \\
 & \textbf{Sqrt-Samp} & \textbf{10.00} & \textbf{22.36} & \textbf{50.34} & 78.24
& 76.12 & {\ul 6.52} & \textbf{59.25} & 90.46 & 90.35 & 0.00 & 51.15 & 48.01 &
0.0 & {\ul 8.06} & 74.14 & 70.95 \\
 & \textbf{CBFocal} & 0.00 & 4.20 & 35.55 & \textbf{87.20} & \textbf{83.62} &
0.00 & 46.51 & {\ul 94.51} & {\ul 94.34} & 0.00 & 57.36 & 53.83 & 0.0 & 4.09 &
\textbf{77.31} & {\ul 73.78} \\
 & \textbf{BAGS} & 1.25 & 7.15 & 32.99 & {\ul 87.01} & {\ul 83.37} &
\textbf{7.25} & 51.45 & \textbf{94.77} & \textbf{94.62} & \textbf{0.26} & {\ul
58.88} & {\ul 55.28} & 0.0 & \textbf{30.05} & 76.00 & \textbf{73.79} \\
 & \textbf{SSB (Ours)} & {\ul 7.50} & {\ul 18.39} & {\ul 42.39} & 86.18 & 83.17
& 5.80
& {\ul 55.96} & 93.81 & 93.67 & {\ul 0.08} & \textbf{59.82} & \textbf{56.15} &
0.0 & 6.75 & 77.01 & 73.62 \\ \hline
\multirow{5}{*}{\textbf{Swin-S}} & \textbf{Baseline} & {\ul 8.75} & 37.12 &
62.99 & \textbf{90.86} & \textbf{88.76} & 10.87 & 60.42 & \textbf{95.09} &
\textbf{94.97} & \textbf{4.17} & {\ul 68.94} & {\ul 64.96} & 0.0 & 7.61 &
\textbf{78.94} & {\ul 75.51} \\
 & \textbf{Sqrt-Samp} & \textbf{16.25} & \textbf{44.84} & \textbf{67.22} & 83.49
& 82.13 & 13.04 & \textbf{64.85} & 89.08 & 88.99 & 1.08 & 64.95 & 61.02 & 0.0 &
{\ul 13.09} & 76.54 & 73.48 \\
 & \textbf{CBFocal} & 0.00 & 25.09 & 56.57 & 90.32 & 87.78 & 4.35 & 53.71 &
94.97 & 94.82 & {\ul 3.78} & \textbf{69.61} & \textbf{65.57} & 0.0 & 5.44 &
78.49 & 74.98 \\
 & \textbf{BAGS} & 6.25 & 21.68 & 52.33 & 90.08 & 87.32 & \textbf{15.22} & 52.32
& 94.81 & 94.67 & 3.67 & 66.01 & 62.18 & 0.0 & \textbf{19.65} & 78.62 &
\textbf{75.78} \\
 & \textbf{SSB (Ours)} & \textbf{16.25} & {\ul 40.18} & {\ul 64.91} & {\ul
90.65} &
{\ul 88.69} & {\ul 11.59} & {\ul 61.94} & {\ul 95.05} & {\ul 94.93} & 1.08 &
67.57 & 63.49 & 0.0 & 7.16 & {\ul 78.86} & 75.41
\\\hline
\end{tabular}}
\end{table}

\end{landscape}


In terms of the macro-averaged F1-score, SSB achieved a score that was higher
than the other methods for most datasets and architectures, as shown in
\autoref{fig:fullimage_f1}. For all the datasets, the Swin-S trained with the
baseline settings attained the best macro F1-score among all architectures, with
Swin-S trained with SSB having the second-best result.
\autoref{fig:wcs_fullimage_f1_resnet} shows the increment/decrement of the
F1-score class-wise over the baseline for ResNet-50 in the WCS dataset after
applying each long-tail method. SSB is the method with the most consistent
increment in terms of F1-score across classes for ResNet-50.
\autoref{fig:wcs_fullimage_f1_swin} shows that the decrease in F1-score for the
Swin-S is generalized over all classes, with SSB being the method that lowers
the F1-score the least.

\subsection{Comparison between methods when using cropped images}

In this experiment, an assessment was made to see whether using an object
detector before classifying an animal can improve the recognition score for the
tail classes. The same settings from the previous experiment were used, except
for using MegaDetectorV4 to generate cropped regions around animals as inputs,
though reverting back to the whole image when the confidence score was less than
0.6. For the pre-cropping experiments, the Snapshot Serengeti dataset was not
used due to the time required to generate bounding boxes with MegaDetectorV4 and
run the experiments with the large quantity of images available.

As shown in \autoref{table:bbox_results}, there was a general accuracy and macro
F1-score increase when the models were trained using cropped images. However,
these results follow the same pattern observed when using the entire image as
the input. Specifically, square-root sampling attained the highest improvements
for the tail classes at the cost of reducing performance for the head classes.
Our proposed SSB was able to improve the tail classes’ accuracy with minimal
accuracy loss for the head classes. In addition, it attained the best macro
F1-score in most of the cases. Lastly, \autoref{fig:image_preds} shows some
sample predictions for pictures randomly picked from the WCS dataset with
bounding boxes generated by MegaDetectorV4.

\begin{landscape}

\begin{table}
\caption[Results attained by different methods and models on different datasets
when using the bounding box with the highest confidence as input.]{Results
attained by different methods and models on different datasets when using the
bounding box with the highest confidence as input. Bounding boxes were generated
by MegaDetectorV4. Bold and underlined values, respectively, indicate the best
and second-best result for each model and dataset combination.}
\label{table:bbox_results}
\setlength\tabcolsep{2pt}
\resizebox{1.5\textwidth}{!}{
\begin{tabular}{ll|llllll|llll|lllll}
\hline
\multicolumn{1}{c}{\textbf{}} & \multicolumn{1}{c|}{\textbf{}} &
\multicolumn{6}{c|}{\textbf{WCS}} & \multicolumn{4}{c|}{\textbf{Caltech CT}} &
\multicolumn{5}{c}{\textbf{Wellington CT}} \\
\multicolumn{1}{c}{\textbf{Model}} & \multicolumn{1}{c|}{\textbf{Training}} &
\multicolumn{1}{c}{\textbf{$Acc_1$}} & \multicolumn{1}{c}{\textbf{$Acc_2$}} &
\multicolumn{1}{c}{\textbf{$Acc_3$}} & \multicolumn{1}{c}{\textbf{$Acc_4$}} &
\multicolumn{1}{c}{\textbf{$Acc_all$}} & \multicolumn{1}{c|}{\textbf{$F_1^m$}} &
\multicolumn{1}{c}{\textbf{$Acc_2$}} & \multicolumn{1}{c}{\textbf{$Acc_4$}} &
\multicolumn{1}{c}{\textbf{$Acc_all$}} & \multicolumn{1}{c|}{\textbf{$F_1^m$}} &
\multicolumn{1}{c}{\textbf{$Acc_2$}} & \multicolumn{1}{c}{\textbf{$Acc_3$}} &
\multicolumn{1}{c}{\textbf{$Acc_4$}} & \multicolumn{1}{c}{\textbf{$Acc_all$}} &
\multicolumn{1}{c}{\textbf{$F_1^m$}} \\ \hline
\multirow{5}{*}{\textbf{ResNet-50}} & \textbf{Baseline} & 0.00 & 21.68 & 52.94 &
90.61 & 87.85 & 0.429 & 0.00 & 60.59 & 56.87 & 0.423 & 0.0 & 8.18 & {\ul 78.52}
& 75.14 & 0.408 \\
 & \textbf{Sqrt-Samp} & {\ul 20.00} & \textbf{43.47} & \textbf{67.19} & 85.79 &
84.28 & 0.473 & \textbf{0.24} & 62.62 & 58.79 & 0.448 & 0.0 & {\ul 22.21} &
75.32 & 72.76 & \textbf{0.431} \\
 & \textbf{CBFocal} & 3.75 & 26.33 & 61.67 & {\ul 90.98} & \textbf{88.65} &
0.468 & 0.00 & 64.32 & 60.37 & 0.458 & 0.0 & 10.77 & 78.47 & 75.21 & {\ul 0.420}
\\
 & \textbf{BAGS} & \textbf{26.25} & 23.16 & 53.33 & \textbf{91.00} & {\ul 88.28}
& 0.414 & {\ul 0.13} & \textbf{67.01} & \textbf{62.91} & {\ul 0.464} & 0.0 &
\textbf{36.98} & 77.92 & \textbf{75.95} & 0.409 \\
 & \textbf{SSB (Ours)} & 17.50 & {\ul 35.41} & {\ul 61.95} & 90.39 & 88.25 &
\textbf{0.501} & 0.11 & {\ul 65.79} & {\ul 61.76} & \textbf{0.467} & 0.0 & 12.23
& \textbf{78.73} & {\ul 75.53} & 0.416 \\ \hline
\multirow{5}{*}{\textbf{MBNetV3}} & \textbf{Baseline} & 0.00 & 26.11 & 55.93 &
\textbf{91.34} & {\ul 88.72} & 0.434 & 0.00 & 64.80 & 60.82 & {\ul 0.471} & 0.0
& 11.89 & {\ul 80.02} & 76.74 & 0.434 \\
 & \textbf{Sqrt-Samp} & \textbf{26.25} & \textbf{42.22} & \textbf{64.41} & 87.30
& 85.57 & {\ul 0.458} & {\ul 0.29} & 57.41 & 53.91 & 0.436 & 0.0 & {\ul 23.89} &
75.67 & 73.17 & \textbf{0.441} \\
 & \textbf{CBFocal} & 1.25 & 19.64 & 56.86 & 91.04 & 88.39 & 0.418 & 0.13 &
61.29 & 57.53 & 0.468 & 0.0 & 11.44 & 79.30 & 76.03 & 0.422 \\
 & \textbf{BAGS} & 13.75 & 22.47 & 45.07 & 90.13 & 87.07 & 0.373 & \textbf{0.42}
& {\ul 65.10} & {\ul 61.12} & 0.461 & 0.0 & \textbf{38.60} & 79.31 &
\textbf{77.34} & 0.403 \\
 & \textbf{SSB (Ours)} & {\ul 21.25} & {\ul 35.19} & {\ul 60.35} & {\ul 91.13} &
\textbf{88.88} & \textbf{0.498} & 0.11 & \textbf{66.59} & \textbf{62.50} &
\textbf{0.489} & 0.0 & 15.98 & \textbf{80.30} & {\ul 77.20} & {\ul 0.438} \\
\hline
\multirow{5}{*}{\textbf{EffV2-B2}} & \textbf{Baseline} & 0.00 & 19.64 & 48.13 &
90.46 & 87.46 & 0.383 & {\ul 0.11} & 61.34 & 57.58 & 0.424 & 0.0 & 8.40 & {\ul
79.61} & 76.18 & 0.427 \\
 & \textbf{Sqrt-Samp} & \textbf{17.50} & \textbf{35.19} & \textbf{65.09} & 88.37
& 86.49 & {\ul 0.469} & 0.08 & 60.33 & 56.63 & 0.437 & 0.0 & {\ul 23.33} & 76.86
& 74.28 & \textbf{0.452} \\
 & \textbf{CBFocal} & 0.00 & 10.33 & 53.58 & \textbf{91.34} & \textbf{88.40} &
0.392 & 0.00 & 64.72 & 60.74 & {\ul 0.468} & 0.0 & 12.53 & 79.39 & 76.17 & 0.430
\\
 & \textbf{BAGS} & 0.00 & 13.17 & 51.87 & {\ul 91.24} & {\ul 88.26} & 0.384 &
\textbf{0.61} & \textbf{68.03} & \textbf{63.89} & \textbf{0.478} & 0.0 &
\textbf{44.22} & 79.15 & \textbf{77.46} & 0.429 \\
 & \textbf{SSB (Ours)} & {\ul 16.25} & {\ul 30.31} & {\ul 60.28} & 90.38 & 88.10
& \textbf{0.474} & 0.08 & {\ul 67.10} & {\ul 62.98} & 0.465 & 0.0 & 12.27 &
\textbf{79.80} & {\ul 76.55} & {\ul 0.439} \\ \hline
\multirow{5}{*}{\textbf{Swin-S}} & \textbf{Baseline} & 16.25 & 53.35 & 78.20 &
{\ul 93.63} & \textbf{92.28} & \textbf{0.637} & {\ul 10.32} & {\ul 70.49} & {\ul
66.80} & \textbf{0.597} & 0.0 & 11.25 & {\ul 80.14} & {\ul 76.82} & 0.503 \\
 & \textbf{Sqrt-Samp} & \textbf{25.00} & \textbf{62.09} & \textbf{79.37} & 90.18
& 89.22 & 0.586 & 2.90 & 67.84 & 63.85 & 0.524 & 0.0 & \textbf{19.05} & 78.42 &
75.56 & {\ul 0.504} \\
 & \textbf{CBFocal} & 8.75 & 39.95 & 75.24 & \textbf{93.66} & {\ul 91.97} & {\ul
0.599} & \textbf{11.40} & 69.41 & 65.84 & {\ul 0.596} & 0.0 & 9.41 & 79.25 &
75.88 & 0.476 \\
 & \textbf{BAGS} & 8.75 & 33.83 & 69.93 & 93.52 & 91.52 & 0.511 & 8.29 &
\textbf{73.86} & \textbf{69.84} & 0.560 & 0.0 & {\ul 18.64} & 79.36 & 76.44 &
0.447 \\
 & \textbf{SSB (Ours)} & {\ul 23.75} & {\ul 58.34} & {\ul 78.55} & 93.53 &
\textbf{92.28} & \textbf{0.637} & 3.27 & 69.52 & 65.45 & 0.576 & 0.0 & 12.34 &
\textbf{80.75} & \textbf{77.46} & \textbf{0.514} \\ \hline
\end{tabular}}
\end{table}

\end{landscape}

\subsection{Comparing SSB to one-stage training strategies}

In this experiment, the performance of single-stage training methods was
evaluated by applying re-weighting (CBFocal) and re-sampling (square-root
sampling). For both methods, cropping was also used before classification, with
bounding boxes being generated by MegaDetectorV4.
\autoref{table:onestage_results} shows the results and compares them to the
baseline training procedure and to our proposed two-stage SSB. The results show
that square-root sampling presents similar behavior as observed in the two-stage
training, i.e., improving the performance of the tail classes, but reducing the
performance of the head classes. On the other hand, CBFocal obtained a great
improvement over its results in two-stage training, and achieved the best
overall accuracy in most cases. However, this improvement occurred mostly in the
head classes. One-stage training methods also fail to improve the classification
performance of the tail classes from the Caltech and Wellington datasets.
Considering the objective of improving accuracy of the tail classes, our
proposed SSB still represents the best trade-off between improving accuracy of
the tail classes and reducing performance in the head classes.


\begin{table}[htp]
\footnotesize
\caption[Comparing the results of one-stage training using square-root
re-sampling and CBFocal loss with the baseline training procedure and our
proposed SSB.]{Comparing the results of one-stage training using square-root
re-sampling and CBFocal loss with the baseline training procedure and our
proposed SSB. Bold and underlined values, respectively, indicate the best and
second-best result for each model and dataset combination.}
\label{table:onestage_results}
\setlength\tabcolsep{2pt}
\resizebox{\textwidth}{!}{
\begin{tabular}{ll|llllll|llll|lllll}
\hline
\multicolumn{1}{c}{\textbf{}} & \multicolumn{1}{c|}{\textbf{}} &
\multicolumn{6}{c|}{\textbf{WCS}} & \multicolumn{4}{c|}{\textbf{Caltech CT}} &
\multicolumn{5}{c}{\textbf{Wellington CT}} \\
\multicolumn{1}{c}{\textbf{Model}} & \multicolumn{1}{c|}{\textbf{Training}} &
\multicolumn{1}{c}{\textbf{$Acc_1$}} & \multicolumn{1}{c}{\textbf{$Acc_2$}} &
\multicolumn{1}{c}{\textbf{$Acc_3$}} & \multicolumn{1}{c}{\textbf{$Acc_4$}} &
\multicolumn{1}{c}{\textbf{$Acc_all$}} & \multicolumn{1}{c|}{\textbf{$F_1^m$}} &
\multicolumn{1}{c}{\textbf{$Acc_2$}} & \multicolumn{1}{c}{\textbf{$Acc_4$}} &
\multicolumn{1}{c}{\textbf{$Acc_all$}} & \multicolumn{1}{c|}{\textbf{$F_1^m$}} &
\multicolumn{1}{c}{\textbf{$Acc_2$}} & \multicolumn{1}{c}{\textbf{$Acc_3$}} &
\multicolumn{1}{c}{\textbf{$Acc_4$}} & \multicolumn{1}{c}{\textbf{$Acc_all$}} &
\multicolumn{1}{c}{\textbf{$F_1^m$}} \\ \hline
\multirow{4}{*}{\textbf{ResNet-50}} & \textbf{Baseline} & 0.00 & 21.68 & 52.94 &
{\ul 90.61} & 87.85 & 0.429 & 0.00 & 60.59 & 56.87 & 0.423 & 0.0 & 8.18 & 78.52
& 75.14 & 0.408 \\
 & \textbf{Sqrt-Samp} & \textbf{20.00} & \textbf{45.86} & \textbf{71.61} & 84.07
& 82.90 & {\ul 0.530} & \textbf{0.90} & 64.87 & 60.94 & {\ul 0.529} & 0.0 & {\ul
13.58} & 77.29 & 74.22 & {\ul 0.458} \\
 & \textbf{CBFocal} & 15.00 & {\ul 35.98} & {\ul 67.58} & \textbf{92.13} &
\textbf{90.15} & \textbf{0.538} & 0.05 & \textbf{71.67} & \textbf{67.27} &
\textbf{0.541} & 0.0 & \textbf{15.00} & \textbf{79.07} & \textbf{75.99} &
\textbf{0.460} \\
 & \textbf{SSB (Ours)} & {\ul 17.50} & 35.41 & 61.95 & 90.39 & {\ul 88.25} &
0.501 & {\ul 0.11} & {\ul 65.79} & {\ul 61.76} & 0.467 & 0.0 & 12.23 & {\ul
78.73} & {\ul 75.53} & 0.416 \\ \hline
\multirow{4}{*}{\textbf{MBNetV3}} & \textbf{Baseline} & 0.00 & 26.11 & 55.93 &
{\ul 91.34} & 88.72 & 0.434 & 0.00 & {\ul 64.80} & {\ul 60.82} & 0.471 & 0.0 &
11.89 & 80.02 & 76.74 & 0.434 \\
 & \textbf{Sqrt-Samp} & \textbf{22.50} & \textbf{48.13} & \textbf{71.75} & 84.85
& 83.67 & \textbf{0.504} & {\ul \textbf{1.08}} & 58.20 & 54.69 & {\ul 0.487} &
0.0 & 14.55 & 76.35 & 73.37 & \textbf{0.459} \\
 & \textbf{CBFocal} & 3.75 & 24.86 & {\ul 61.99} & \textbf{91.69} &
\textbf{89.31} & 0.463 & 0.08 & 64.23 & 60.29 & 0.485 & 0.0 & {\ul 15.64} & {\ul
80.11} & {\ul 77.00} & {\ul 0.452} \\
 & \textbf{SSB (Ours)} & {\ul 21.25} & {\ul 35.19} & 60.35 & 91.13 & {\ul 88.88}
& {\ul 0.498} & {\ul 0.11} & \textbf{66.59} & \textbf{62.50} & \textbf{0.489} &
0.0 & \textbf{15.98} & \textbf{80.30} & \textbf{77.20} & 0.438 \\ \hline
\multirow{4}{*}{\textbf{EffV2-B2}} & \textbf{Baseline} & 0.00 & 19.64 & 48.13 &
{\ul 90.46} & 87.46 & 0.383 & {\ul 0.11} & 61.34 & 57.58 & 0.424 & 0.0 & 8.40 &
79.61 & 76.18 & 0.427 \\
 & \textbf{Sqrt-Samp} & \textbf{26.25} & \textbf{56.64} & \textbf{78.09} & 85.66
& 84.84 & \textbf{0.565} & \textbf{2.61} & 61.37 & 57.76 & {\ul 0.500} & 0.0 &
{\ul 16.54} & 77.07 & 74.15 & {\ul 0.460} \\
 & \textbf{CBFocal} & 0.00 & 13.85 & 60.24 & \textbf{92.37} & \textbf{89.72} &
0.410 & 0.00 & \textbf{71.39} & \textbf{67.00} & \textbf{0.529} & 0.0 &
\textbf{18.60} & \textbf{80.75} & \textbf{77.76} & \textbf{0.471} \\
 & \textbf{SSB (Ours)} & {\ul 16.25} & {\ul 30.31} & {\ul 60.28} & 90.38 & {\ul
88.10} & {\ul 0.474} & 0.08 & {\ul 67.10} & {\ul 62.98} & 0.465 & 0.0 & 12.27 &
{\ul 79.80} & {\ul 76.55} & 0.439 \\ \hline
\multirow{4}{*}{\textbf{Swin-S}} & \textbf{Baseline} & 16.25 & 53.35 & 78.20 &
{\ul 93.63} & {\ul 92.28} & {\ul 0.637} & \textbf{10.32} & {\ul 70.49} & {\ul
66.80} & {\ul 0.597} & 0.0 & 11.25 & 80.14 & 76.82 & 0.503 \\
 & \textbf{Sqrt-Samp} & \textbf{23.75} & 51.87 & \textbf{78.62} & 92.03 & 90.78
& 0.634 & 2.16 & 67.71 & 63.68 & 0.574 & 0.0 & 11.29 & 79.36 & 76.08 & 0.479 \\
 & \textbf{CBFocal} & {\ul 21.25} & {\ul 54.03} & 78.41 & \textbf{93.97} &
\textbf{92.62} & \textbf{0.642} & {\ul 6.76} & \textbf{71.62} & \textbf{67.63} &
\textbf{0.606} & 0.0 & \textbf{13.43} & {\ul 80.56} & {\ul 77.33} &
\textbf{0.524} \\
 & \textbf{SSB (Ours)} & \textbf{23.75} & \textbf{58.34} & {\ul 78.55} & 93.53 &
{\ul 92.28} & {\ul 0.637} & 3.27 & 69.52 & 65.45 & 0.576 & 0.0 & {\ul 12.34} &
\textbf{80.75} & \textbf{77.46} & {\ul 0.514} \\ \hline
\end{tabular}}
\end{table}

\section{Discussion}

The long-tail distribution is a characteristic of real-world tasks, such as
visual recognition of species in camera-trap images, in which some classes
concentrate the majority of available instances and a large number of classes
have just a small quantity of samples. However, as indicated by
\citet{schneider2020three}, computer vision models generally often perform
poorly on recognizing species in camera-trap images with fewer training
instances (tail classes). In this Chapter, a simple and effective framework
called Square-root Sampling Branch (SSB) is proposed, which can improve the
performance in tail classes with a minimal percentage of accuracy loss in the
head classes. The proposed framework was systematically evaluated against
state-of-the-art methods for handling the long-tail distribution issue in four
different camera-trap datasets and four families of computer vision models. Our
experiments show that our proposed framework presents competitive results with
the methods in the literature, achieving the best trade-off between increasing
the performance of the tail classes and reducing the head classes’ accuracy in
most cases.

Comparing our improved baseline training procedure with the literature,
ResNet-50 achieved an overall accuracy of 92.79\% for the Snapshot Serengeti
dataset using the full image, while the same architecture achieved an accuracy
of 93.6\% in the work conducted by \citet{norouzzadeh2018automatically} and
approximately 85\% in that of \citet{villa2017towards}. However, while these
works use burst or random split, in our methodology we split the training/test
partitions based on locations, which tends to achieve a significantly lower
accuracy, as demonstrated for other datasets by \citet{beery2018recognition},
\citet{schneider2020three}, and \citet{tabak2019machine}, in which the
performance of models on untrained locations drops from 15\% to 25\% in accuracy
in terms of absolute values. Furthermore, the categories and test sets used are
different among the works, thus making the results not directly comparable. 

The
same occurs when comparing our results for the Caltech CT dataset with the
results from \citet{beery2018recognition}. The lack of well-defined benchmarks
for visual recognition of animal species in camera-trap images makes it
difficult to compare the results of different methods across works. Although
most of the applications have their specificities and requirements regarding the
data preparation or even the target dataset, it would be important to apply them
using a standardized benchmark as well. However, designing such a benchmark
dataset is challenging due to the wide variety of factors that need to be
considered, including taxonomic coverage, biogeography, habitat and climate
conditions, as well as the tasks to be included, such as species classification,
long-tail visual recognition, animal detection, counting, and behavior
identification. Therefore, it is recommended that the research community work
collaboratively towards developing standardized benchmarks for automatic
information extraction from camera-trap images, which would facilitate
comparisons between different methods and ultimately lead to advancements in
this field.

In general, the method that most improved the performance for the tail classes
over the baseline was the square-root sampling. However, this method strongly
degrades the head classes’ accuracy. This behavior may be due to the high
imbalance of the datasets, leading the square-root sampling to cause a high
undersampling of the majority classes. Our proposed SSB can improve the accuracy
for the tail classes more than CBFocal and BAGS for most of the cases, but,
unlike square-root sampling, the degradation caused to the head classes is
minimal, representing the best trade-off among all the evaluated methods. SSB
also represents the best performance in terms of macro F1-score for most of the
cases when compared with the two-stage methods. On the other hand, our
experiments using one training stage show that CBFocal achieved the best overall
accuracy for most of the cases, primarily because of its superior performance in
the head classes. 

This CBFocal performance improvement may indicate that
pretraining using softmax cross-entropy does not generate a good representation
for fine-tuning using Focal loss, or that the learning of representation and
classification is not well decoupled when using this loss function. However, our
experiments were not designed to assess the performance based on representation
(feature extractor) training strategies. Furthermore, \citet{Kang2020Decoupling}
do not evaluate the performance of re-weighting methods when proposing the
two-stage training approach. Although this approach was not explored in our
experiments, a potential alternative could be to combine the outcomes of
square-root resampling and CBFocal, which respectively performed better on the
tail and head classes. However, this approach would require running inference
using two models, effectively doubling the computational cost. In contrast, our
proposed SSB only needs to be run once and is still competitive.

Looking into the model’s performance, Swin Transformer obtained the best results
considering all the datasets, which was expected since this model has more
parameters than the other selected models. It is worth noting that the Swin-S
that was trained using the baseline settings attained an excellent performance
even for the tail classes, reaching the best mean F1-score in most of the
datasets. This superior performance of vision transformers in long-tail
classification corroborates with what has been stated in the literature, as in
the classification of fungi \citep{picek2022danish}, for instance.

Regarding the question of whether using cropping before classification can
improve model performance, when comparing the results of the first and the
second series of experiments, it is possible to note a significant increase in
the overall accuracy when applying cropping. This improvement is more relevant
for the WCS and Caltech datasets, in which, for some cases, there was an
improvement of 10\% and 15\%, respectively. This performance improvement is
consistent with the results in the literature, for example, in
\citep{beery2018recognition}, there was a 20\% increase in the accuracy when
cropping was used before classification. The use of an object detector as a
preprocessing step before species identification is also part of the methodology
of \citet{schneider2018deep} and \citet{norouzzadeh2021deep}; however, its
impact on the performance is not evaluated. Looking at the binned accuracy
results, a performance increase is observed for all bins of the classes, except
for the tail classes ($Bin_2$) for the Caltech CT and Wellington CT datasets, in
which all models and methods fail to learn them. The lowest performance
increases due to the use of cropping before classification were observed for the
Swin-S model, with an improvement varying from 1\% to 4\%. It is worth noting
that the accuracies for Swin-S using the full images are already very high when
compared to the other models. This can be explained by the vision transformer’s
attention mechanism that can allow it to focus on the animal even when using the
full image.

Moreover, our experiments did not explain why the models were unable to learn to
recognize the tail classes in the smaller datasets. The Caltech and Wellington
datasets, with only 19 and 14 classes respectively, had a limited number of tail
classes (4 and 1, respectively). It is possible that these tail classes were
ignored due to the relatively low number of training samples available. To
investigate this further, future studies could use the WCS dataset and
progressively remove classes to evaluate the impact on model performance. This
would also provide insights into the minimum number of training images per
class needed to see an improvement in performance when using the evaluated
methods.

One limitation to our study is that only the WCS dataset has classes belonging
to all bins of classes in regards to the number of training instances. Another
limitation is that, during our experiments, there was an effort to keep our
settings as similar as possible across the methods and models in order to make a
fair comparison; therefore, our hyperparameters are suboptimal and better tuning
would lead to improvements.

\section{Final Remarks}

Although rare or endangered species may be the classes of interest for projects
using camera traps for monitoring wildlife, they are often neglected when
developing deep-learning models for automatic species recognition, mostly
because of the difficulty involved in learning classes with few training
instances from long-tail datasets. In this chapter, a simple and effective
framework called Square-root Sampling Branch (SSB) is proposed, which uses a
multi-branch approach to improve the performance recognition of tail classes of
long-tail camera-trap datasets with minimal loss in the performance recognition
of the head classes. Our proposed method was systematically evaluated against
square-root sampling, CBFocal loss, and Balanced Group Softmax (BAGS)
considering four camera-trap datasets: Snapshot Serengeti, WCS, Caltech CT, and
Wellington CT. Our results showed that our proposed approach achieved the best
trade-off between improving accuracy for the tail classes and decreasing the
accuracy of the head classes, and is competitive with the state-of-the-art
methods. It was also found that Swin-S performed better for most of the datasets
even without applying any additional method for handling imbalance, and achieved
an overall accuracy of 88.76\% for the WCS dataset and 94.97\% for the Snapshot
Serengeti, when considering location-based splitting.

Despite the fact that our proposed framework improves the recognition of tail
classes, our experiments clearly show that it is necessary to develop more
sophisticated methods to handle the long-tail distribution in the recognition of
species in camera-trap images. Specifically for the tail classes ($Bin_1$ and
$Bin_2$), for which the results were extremely low, and especially for the small
datasets, one could evaluate few-shot learning methods. Another issue that was
discussed is the lack of standardized benchmarks for evaluating the performance
of long-tail visual recognition of animal species when comparing new methods. It
is hoped that our well-defined methodology for dataset preprocessing and
training procedure helps establish good baselines for other researchers. In
future works, other methods can be explored such as progressive-balanced
sampling, Ally Complementary Experts, and Rebalanced Mixup (Remix). Additional
training tricks, such as label smoothing and mixup, can also be performed, but
they must be used with caution because one trick may affect the other.

The next chapter explores the task of counting the number of individuals of each 
species present in a sequence of camera-trap images. For this task, a good model 
for species identification that can handle long-tail distributions is essential, 
since the task considers not only the total number of individuals but also whether 
the count is correct for each species.
